\clearpage
\onecolumn
\thispagestyle{empty}
\vspace*{\fill}
\begin{center}
	{\color{radblue}\rule{\linewidth}{2pt}}\par
	\vspace{6mm}
	{\Huge\bfseries\color{radblue} Mamografía}\par
	\vspace{6mm}
	{\color{radblue}\rule{\linewidth}{2pt}}
\end{center}
\vspace*{\fill}
\clearpage
\twocolumn
\pagestyle{contenido}

%%%%%%%%%%%%%%%%%%%%%%%%%%%%%%%%%%%%%%%%%%%%%%%%%%%%%%%%%%%
%  SECCIÓN: Mamografía
%  Archivo: secciones/mamografia.tex
%  Sin preámbulo ni \begin{document} — lo maneja main.tex
	%%%%%%%%%%%%%%%%%%%%%%%%%%%%%%%%%%%%%%%%%%%%%%%%%%%%%%%%%%%
	
	\section{Mamografía}
	
	\subsection{Introducción a la Mamografía}
	
	\subsubsection*{1. Definición y Rol Clínico}
	
	\textbf{Técnica:} Principal método de estudio de la mama mediante imagen. Técnica de elección para la detección precoz del cáncer de mama.
	
	\textbf{Alcance diagnóstico:} La mamografía \textbf{no diagnostica} el cáncer de mama. El diagnóstico definitivo requiere confirmación anatomopatológica mediante biopsia del tejido de la lesión.
	
	\textbf{Comunicación con el paciente:} Siempre transmitir el mensaje con claridad, reducir el miedo y promover activamente la prevención.
	
	\subsubsection*{2. Anatomía de la Mama}
	
	\paragraph{Estructura general}
	Glándula lobulada de forma hemisférica, constituida por estroma fibroso con tejido adiposo. Se localiza en la aponeurosis superficial de la región anterolateral del tórax, tanto femenino como masculino. Presenta una prolongación axilar denominada \textbf{cola de Spence}, zona de ganglios linfáticos con relevancia oncológica.
	
	
	\paragraph{Organización glandular}
	La mama femenina adulta consta de \textbf{15 a 20 lóbulos}, cada uno subdividido en lobulillos individuales. Estos contienen los elementos glandulares denominados \textbf{acinos} o \textbf{ULDT} (Unidades Lobulillares Ductales Terminales), estructuras de importancia clínica por ser origen frecuente de carcinomas. Los conductos desembocan en el pezón.
	
	\paragraph{Ligamentos de Cooper}
	Las mamas están sujetas por los \textbf{ligamentos de Cooper}. La firmeza o flacidez de la mama depende del estado de estos ligamentos y del contenido relativo de tejido adiposo.
	
	\subsubsection*{3. Clasificación por Tipo de Mama}
	
	\paragraph{Mama A — Adiposa}
	Totalmente grasa. Radiolúcida (hipodensa). Presente en pacientes post-menopausia, mamas infantiles y masculinas. Requiere menor compresión.
	
	\paragraph{Mama B — Fibrosa / Densa}
	Predominio de tejido fibroso. Mayor densidad radiológica.
	
	\paragraph{Mama C — Heterogéneamente Densa}
	Densidad heterogénea. Fibroglandular. Requiere estudio complementario con ecografía. Mayor propensión a desarrollar cáncer.
	
	\paragraph{Mama D — Muy Densa / Fibroglandular}
	Hiperdensa, de mayor complejidad diagnóstica. Presente en pacientes post-pubertad hasta los 30 años, embarazadas, período de lactancia y por condicionamiento genético. La compresión puede resultar molesta. Requiere ecografía complementaria obligatoria. Mayor propensión a desarrollar cáncer.
	
	\subsubsection*{4. Signos Clínicos a Identificar}
	
	Al realizar el estudio, registrar y describir toda señal visible. Los principales signos a identificar son:
	
	\begin{itemize}
		\item Engrosamiento de la piel
		\item Secreción del pezón (acuosa, láctea o hemática), unilateral o bilateral
		\item Dilatación venosa superficial
		\item Hundimiento o retracción del pezón
		\item Protuberancia o masa palpable
		\item Cambio de forma o tamaño
		\item Costras en el pezón
		\item Piel de naranja (signo de edema dérmico con aspecto punteado)
	\end{itemize}
	

	
	\subsubsection*{5. Sistema BI-RADS}
	
	\textbf{Definición:} Sistema estandarizado de categorización de hallazgos mamográficos. Escala del \textbf{0 al 6}.
	
	\textbf{Categorías:}
	
	\begin{itemize}
		\item \textbf{0} — Estudio insuficiente. Se recomienda correlación con otro método de imagen.
		\item \textbf{1 y 2} — Benigno. Sin hallazgos significativos o hallazgo con características benignas definidas.
		\item \textbf{3} — Intermedio. Lesión con apariencia de benignidad; requiere control evolutivo cada 6 meses.
		\item \textbf{4} — Sospecha moderada. Subcategorías A, B y C según grado de sospecha.
		\item \textbf{5} — Altamente sospechoso de malignidad.
		\item \textbf{6} — Malignidad confirmada (diagnóstico histológico previo).
	\end{itemize}
	
	\subsubsection*{6. Consideraciones Clínicas Relevantes}
	
	\textbf{Límite diagnóstico:} La mamografía detecta lesiones pero no las diagnostica. Toda categoría BI-RADS 4 o superior requiere derivación para biopsia.
	
	\textbf{Mamas densas (C y D):} Reducen la sensibilidad de la mamografía por enmascaramiento de lesiones. La ecografía complementaria es obligatoria en estos casos.
	
	\textbf{Cola de Spence:} Zona axilar de la mama con presencia de ganglios linfáticos. Área de relevancia oncológica; puede ser asiento de metástasis ganglionares.
	
	\textbf{Prevención:} Promover activamente el control mamográfico periódico. La detección precoz modifica significativamente el pronóstico.
	
\clearpage
\onecolumn

\begin{figure}[]
	\centering
	\includegraphics[width=0.5\linewidth]{imagenes/mamo1}
	\caption[]{Tipos de mama}
	\label{fig:mamo1}
\end{figure}

	\begin{figure}[]
	\centering
	\includegraphics[width=0.40\linewidth]{imagenes/mamo2}
	\caption{Signos clínicos}
	\label{fig:mamo2}
\end{figure}

\twocolumn
\clearpage
	
	%----------------------------------------------------------
	%  Para agregar otro enfoque duplicá el bloque
	%  desde \subsection hasta \textbf{Trucos}
	%----------------------------------------------------------
	
	
%%%%%%%%%%%%%%%%%%%%%%%%%%%%%%%%%%%%%%%%%%%%%%%%%%%%%%%%%%%
%  SECCIÓN: Mamografía
%  Archivo: secciones/mamografia.tex
%  Sin preámbulo ni \begin{document} — lo maneja main.tex
	%%%%%%%%%%%%%%%%%%%%%%%%%%%%%%%%%%%%%%%%%%%%%%%%%%%%%%%%%%%
	
	\section{Mamografia}
	
	\subsection{Proyección Cráneo-Caudal (CC) Der / Izq}
	
	\subsubsection*{1. Identificación}
	
	\textbf{Tipo:} Proyección básica de mamografía. Forma par con la proyección oblicua mediolateral (OML).
	
	\textbf{Indicaciones clínicas principales:}
	\begin{itemize}
		\item Detección y evaluación de calcificaciones, quistes y carcinomas
		\item Identificación de otras anomalías o cambios del tejido mamario
		\item Estudio de nódulos palpables
		\item Estudio bilateral de rutina (screening o diagnóstico)
	\end{itemize}
	
	\subsubsection*{2. Técnica Base}
	
	\textbf{Receptor de imagen:} 18×24 cm o 24×30 cm según el tamaño de la mama. Orientación \textbf{transversal}.
	
	\textbf{DFRI:} \textit{[dato no suministrado]}.
	
	\textbf{Técnica:} Ultra contrastada. La mama está compuesta por tejidos con bajo contraste intrínseco (grasa y glándula son radiológicamente homogéneos); la técnica ultra contrastada --- mediante kVp bajo sumado a la compresión --- genera una escala de grises corta que permite diferenciarlos entre sí.
	
	\textbf{Factores técnicos:} kVp \textbf{\textit{[dato no suministrado]}} --- mAs \textbf{\textit{[dato no suministrado]}}. Espesor medido con la mama comprimida sobre el receptor.
	
	\textbf{Bucky:} Equipo de mamografía con compresor dedicado; grilla específica de mamografía.
	
	\subsubsection*{3. Posicionamiento}
	
	\paragraph{Paciente}
	De pie o sentado, con ambos pies orientados frente al equipo. Brazos relajados a los lados con descenso activo del hombro ipsilateral. Cabeza girada hacia el lado opuesto al seno a estudiar.
	
	\paragraph{Región de interés}
	La mama se coloca sobre el receptor, perpendicular a la caja torácica. Se rota levemente el pezón hasta dejarlo de perfil. Se realiza tracción del tejido mamario hacia el receptor, alejándolo de la pared torácica. El pliegue inframamario se eleva al nivel del borde del receptor.
	
	\paragraph{Rayo Central}
	Perpendicular a la base de la mama (plano horizontal). Punto de entrada: \textbf{base de la mama, a nivel del pliegue inframamario elevado}.
	
	\paragraph{Receptor de imagen}
	Elevado para apoyar el pliegue inframamario. Centrado de forma que incluya todo el parénquima mamario, la grasa retromamaria y, de ser posible, el músculo pectoral mayor en el borde posterior. El pedal del compresor se acciona una vez realizada la tracción y correcta posición de la mama.
	
	\paragraph{Respiración}
	Suspendida (paciente contiene la respiración durante la exposición). Reduce el movimiento y evita artefactos de borrosidad.
	
	\subsubsection*{4. Anatomía Visible}
	
	Parénquima mamario completo en proyección cráneo-caudal. Grasa retromamaria en el borde posterior. Músculo pectoral mayor visible en el margen posterior cuando la tracción y la elevación del pliegue son adecuadas. Pezón de perfil y centrado en el borde anterior. Piel y tejido subcutáneo en los bordes superior e inferior.
	
	\subsubsection*{5. Criterios de Evaluación}
	
	\textbf{Posicionamiento correcto:} Ambas proyecciones CC simétricas entre sí. Pezón de perfil y centrado. Bordes superior e inferior (pliegue inframamario) contenidos dentro del campo. Músculo pectoral mayor visible en el borde posterior. La medida de la PNL en la CC no debe diferir en más de 1 cm respecto a la PNL de la proyección OML del mismo lado.
	
	\textbf{Ausencia de rotación:} Pezón centrado y de perfil; si no está de perfil indica falta de rotación, corregible con rotación externa o interna leve de la mama. Densidad glandular simétrica entre ambas CC; asimetría puede indicar error de posicionamiento o variación anatómica del paciente.
	
	\textbf{Penetración y contraste:} Escala de grises corta que permita diferenciar tejido glandular de tejido graso. Técnica ultra contrastada verificable por la correcta separación tonal entre componentes mamarios.
	
	\subsubsection*{6. Puntos Críticos}
	
	\textbf{Pezón no de perfil:} La rotación insuficiente impide evaluar correctamente el tejido retroareolar; corregir con rotación externa o interna leve antes de comprimir.
	
	\textbf{Tracción insuficiente o elevación incorrecta del pliegue:} El músculo pectoral no se visualiza en el borde posterior y se excluye tejido mamario profundo; realizar nueva tracción activa y elevar correctamente el pliegue inframamario antes de comprimir.
	
	\textbf{Adaptaciones:} \textit{[dato no suministrado]}.
	
	\textbf{Trucos:} Verificar la PNL de la CC contra la PNL de la OML del mismo lado: una diferencia mayor a 1 cm indica pérdida de tejido profundo en alguna de las dos proyecciones y obliga a repetir. La asimetría de densidad entre ambas CC puede deberse a la anatomía del paciente o a un error de posicionamiento encubierto.
	
	\clearpage
	\onecolumn
	
	\begin{figure}
		\centering
		\includegraphics[width=1\linewidth]{imagenes/mamo4}
		\caption{CC: craneo caudal + I o D, MLO: medio lateral oblicua + I o D}
		\label{fig:mamo4}
	\end{figure}
	
	\twocolumn
	\clearpage
	
	%----------------------------------------------------------
	%  Para agregar otro enfoque duplicá el bloque
	%  desde \subsection hasta \textbf{Trucos}
	%----------------------------------------------------------
	
%%%%%%%%%%%%%%%%%%%%%%%%%%%%%%%%%%%%%%%%%%%%%%%%%%%%%%%%%%%
%  SECCIÓN: Mamografía
%  Archivo: secciones/mamografia.tex
%  Sin preámbulo ni \begin{document} — lo maneja main.tex
	%%%%%%%%%%%%%%%%%%%%%%%%%%%%%%%%%%%%%%%%%%%%%%%%%%%%%%%%%%%
	
	\section{Mamografia}
	
	\subsection{Proyección Mediolateral Oblicua (MLO)}
	
	\subsubsection*{1. Identificación}
	
	\textbf{Tipo:} Proyección básica de screening mamográfico. Forma par con la proyección craneocaudal (CC).
	
	\textbf{Indicaciones clínicas principales:}
	\begin{itemize}
		\item Screening mamográfico de rutina
		\item Dolor localizado en mama
		\item Sospecha de lesiones en cuadrantes externos y cara lateral profunda del tejido mamario
	\end{itemize}
	
	\subsubsection*{2. Técnica Base}
	
	\textbf{Receptor de imagen:} 18×24 cm o 24×30 cm. Orientación \textbf{transversal (horizontal)} según el biotipo de la paciente.
	
	\textbf{DFRI:} \textit{[dato no suministrado]}.
	
	\textbf{Técnica:} Bajo kVp para maximizar la visualización de la cola de Spence y los ganglios axilares. Los tejidos mamarios (grasa y glándula) son anatómicamente distintos pero radiológicamente homogéneos; el kVp bajo combinado con la compresión permite diferenciarlos, logrando una imagen ultra contrastada con escala larga pero de corto rango efectivo.
	
	\textbf{Factores técnicos:} kVp \textbf{bajo} --- mAs \textit{[dato no suministrado]}. \textit{[Método de medición del espesor no suministrado]}.
	
	\textbf{Bucky:} \textit{[dato no suministrado]}.
	
	\subsubsection*{3. Posicionamiento}
	
	\paragraph{Paciente}
	De pie; si no es posible, sentada. Cuerpo y cadera oblicuados hacia el equipo. El brazo del lado a estudiar se eleva sobre el soporte del equipo o sobre la cabeza. Hombro, brazo y cuerpo deben estar completamente relajados; se indica a la paciente un movimiento de rotación del húmero para favorecer la relajación del músculo pectoral mayor.
	
	\paragraph{Región de interés}
	Tubo y receptor rotados a 45° (rango aceptable: 30°--60°) a la altura del hueco axilar. La mama se tracciona en sentido anteromedial, alejándola de la pared torácica, y se coloca sobre el soporte del equipo para quedar bien perfilada. El pezón debe quedar de perfil. La compresión debe lograr grosor uniforme desplazando la mama hacia arriba y afuera.
	
	\paragraph{Rayo Central}
	Perpendicular al plano del receptor, sin angulación (RC inmóvil, 0°). Punto de entrada: \textbf{base de la mama al borde de la pared torácica}.
	
	\paragraph{Receptor de imagen}
	Centrado en la base de la mama al borde de la pared torácica, posicionado a 45° a la altura del hueco axilar. Debe incluir toda la mama desde el músculo pectoral mayor hasta el pezón, la cola de Spence y los ganglios axilares.
	
	\paragraph{Respiración}
	Apnea durante la exposición. Evita el movimiento y la borrosidad de las estructuras mamarias.
	
	\subsubsection*{4. Anatomía Visible}
	
	Toda la mama incluyendo tejido glandular y grasa subcutánea. Músculo pectoral mayor con convexidad anterior visible, indicando relajación adecuada del hombro y la axila. Cola de Spence y ganglios axilares. Pezón de perfil. Pliegue inframamario abierto y despejado, sin superposición del abdomen.
	
	\subsubsection*{5. Criterios de Evaluación}
	
	\textbf{Posicionamiento correcto:} Toda la mama incluida desde el músculo pectoral mayor hasta el pezón; cola de Spence y ganglios axilares visibles; pliegue inframamario abierto y libre de superposición abdominal.
	
	\textbf{Ausencia de rotación:} Pezón de perfil; mama desplazada hacia arriba y afuera con grosor uniforme producto de compresión óptima; músculo pectoral mayor con borde convexo anterior.
	
	\textbf{Penetración y contraste:} Imagen ultra contrastada con diferenciación visible entre tejido glandular y grasa, lograda mediante kVp bajo y compresión adecuada.
	
	\subsubsection*{6. Puntos Críticos}
	
	\textbf{Hombro y brazo tensionados:} El músculo pectoral mayor aparece recto o cóncavo en lugar de convexo, reduciendo el tejido mamario incluido. Corrección: indicar rotación activa del húmero antes de la compresión.
	
	\textbf{Mama caída o mal traccionada:} Grosor no uniforme y pliegue inframamario superpuesto al abdomen. Corrección: traccionar la mama en sentido anteromedial antes de aplicar la compresión.
	
	\textbf{Adaptaciones:} Si la paciente no puede estar de pie, realizar la proyección sentada manteniendo la angulación del tubo a 45°. Para movilidad limitada del brazo, ajustar el ángulo del tubo dentro del rango 30°--60°.
	
	\textbf{Trucos:} Colocar el delantal de protección abdominal antes de la exposición. Verificar el pezón de perfil y la convexidad del pectoral como indicadores rápidos de posicionamiento correcto antes de comprimir definitivamente.
	
	\clearpage
	\onecolumn
	\begin{figure}
		\centering
		\includegraphics[width=0.7\linewidth]{imagenes/mamo5}
		\caption{Medio Lateral Oblicua}
		\label{fig:mamo5}
	\end{figure}
	\twocolumn
	\clearpage
	%----------------------------------------------------------
	%  Para agregar otro enfoque duplicá el bloque
	%  desde \subsection hasta \textbf{Trucos}
	%----------------------------------------------------------
	
	
	%%%%%%%%%%%%%%%%%%%%%%%%%%%%%%%%%%%%%%%%%%%%%%%%%%%%%%%%%%%
	%  SECCIÓN: [Nombre de la Sección]
	%  Archivo: implante_desplazado_eklund.tex
	%  Sin preámbulo ni \begin{document} — lo maneja main.tex
		%%%%%%%%%%%%%%%%%%%%%%%%%%%%%%%%%%%%%%%%%%%%%%%%%%%%%%%%%%%
		
		\section{Mamografia}
		
		\subsection{Implante Desplazado (ID) — Método de Eklund DC}
		
		\subsubsection*{1. Identificación}
		
		\textbf{Indicaciones clínicas principales:}
		\begin{itemize}
			\item Pacientes con implantes mamarios
			\item Aumento mamario
			\item Fuga intra o extracapsular del implante
		\end{itemize}
		
		\subsubsection*{2. Técnica Base}
		
		\textbf{Receptor de imagen:} 18x24 cm o 24x30 cm.
		
		\textbf{Técnica:} Manual o semiautomática (el equipo no detecta el implante).
		
		\subsubsection*{3. Posicionamiento}
		
		\paragraph{Paciente}
		De pie o sentada en taburete, de cara al portacartuchos. Se realiza tracción hacia atrás contra la pared torácica.
		
		\paragraph{Región de interés}
		El operador se coloca de pie en el lado interno de la mama a explorar. Se levanta el pliegue inframamario. Una vez localizado el borde anterior del implante, se tira con suavidad del tejido mamario anterior hacia el portacartucho. Se debe conseguir una compresión adicional de 2--5 cm con el implante desplazado. Se le pide al paciente que indique si la compresión le resulta incómoda.
		
		\paragraph{Respiración}
		Suspendida (contenida). Se solicita al paciente que aguante la respiración una vez conseguida la compresión completa, momento en el que se realiza la exposición.
		
		\subsubsection*{4. Anatomía Visible}
		
		Tejido mamario anterior y central sin superposición con el implante. Tejido glandular anterior al implante desplazado.
		
		\subsubsection*{5. Criterios de Evaluación}
		
		\textbf{Posicionamiento correcto:} Implante desplazado posteriormente contra la pared torácica. Tejido mamario anterior y central visualizado sin ninguna superposición del implante.
	
	\onecolumn
\begin{figure}
	\centering
	\includegraphics[width=0.7\linewidth]{imagenes/mamo6}
	\caption{}
	\label{fig:mamo6}
\end{figure}

\begin{figure}
	\centering
	\includegraphics[width=0.7\linewidth]{imagenes/mamo7}
	\caption{}
	\label{fig:mamo7}
\end{figure}

\twocolumn
\clearpage
		
		%----------------------------------------------------------
		%  Para agregar otro enfoque duplicá el bloque
		%  desde \subsection hasta \textbf{Trucos}
		%----------------------------------------------------------